\documentclass[11pt,a4paper]{article}
\usepackage[utf8]{inputenc}
\usepackage[T1]{fontenc}
\usepackage[spanish,es-noquoting]{babel}
\usepackage{lmodern}
\usepackage{microtype}
\usepackage[margin=2.8cm]{geometry}
\usepackage{enumitem}
\usepackage[hidelinks]{hyperref}

\title{Dirk Obbink, Safo y los papiros de Oxirrinco}
\author{}
\date{}

\begin{document}
\maketitle

\section*{El protagonista}

Dirk Obbink es un papirólogo estadounidense. Fue profesor de Papirología y Literatura Griega en Oxford (Christ Church) y recibió una beca MacArthur en 2001. Durante años fue uno de los editores principales de los papiros de Oxirrinco, la gran colección de la Egypt Exploration Society (EES) depositada en la biblioteca Sackler de Oxford. Era una de las figuras más prestigiosas del campo.

\section*{El «nuevo Safo» (2014)}

En 2014 Obbink anunció un hallazgo espectacular: un papiro con el llamado \textbf{«Poema de los hermanos»}, en el que la voz poética habla de sus hermanos Caraxo y Lárico. El mismo papiro contenía partes del \textbf{«Poema de Cipris»}. Fue el descubrimiento más importante de Safo en décadas, y la prensa de todo el mundo lo recogió.

El problema era la \textbf{procedencia}. Obbink dio varias versiones, que no coincidían entre sí:
\begin{itemize}[nosep]
  \item que el papiro pertenecía a un coleccionista anónimo de Londres;
  \item que se había comprado en una venta privada de Christie's en 2011;
  \item que procedía de un cartonaje de momia (papel reciclado para envolver momias) que se había disuelto para extraer los papiros;
  \item que tenía un historial documentado que lo situaba fuera de Egipto antes de 1970. Esto es importante, porque con ese historial la exportación habría sido legal según la Convención de la UNESCO.
\end{itemize}

Investigadores como Brent Nongbri y Roberta Mazza señalaron las contradicciones. En 2020, C.~Michael Sampson publicó un estudio en el \emph{Bulletin of the American Society of Papyrologists} que desmontó el relato de Christie's y mostró que la historia de la procedencia no se sostenía. Además, otros fragmentos del mismo rollo de Safo habían acabado en la \textbf{colección Green}, de la familia propietaria de Hobby Lobby. Esa conexión llevó directamente al escándalo mayor.

\section*{El «Marcos del siglo I» y Hobby Lobby}

\begin{description}[nosep]
  \item[2012:] Corrió el rumor, difundido entre otros por Scott Carroll (asesor de la colección Green), de que existía un fragmento del Evangelio de Marcos del siglo~I. Sería el manuscrito más antiguo del Nuevo Testamento.
  \item[2013:] Más tarde apareció un contrato, firmado aparentemente por Obbink, para vender ese fragmento y otros a la familia Green, que estaba reuniendo fondos para el \textbf{Museo de la Biblia} de Washington.
  \item[2018:] La EES publicó el fragmento (P.Oxy.\ LXXXIII 5345). Lo dató en los siglos II--III, no en el I, y afirmó que siempre había sido suyo y nunca se había vendido.
  \item[2019:] La EES anunció que al menos 13 papiros de su colección, la mayoría bíblicos, habían acabado en el Museo de la Biblia sin autorización. Las ventas se atribuían a Obbink, y el museo los devolvió. Las revisiones posteriores elevaron a más de un centenar los papiros que faltaban o tenían catalogación alterada. Según la EES, se habían eliminado fotos y registros.
\end{description}

\section*{Consecuencias}

\begin{itemize}[nosep]
  \item La EES le apartó del acceso a la colección en 2016 y después le retiró como editor.
  \item Oxford le suspendió de sus funciones.
  \item En 2020 la policía del Valle del Támesis le detuvo por sospecha de robo y fraude, y quedó en libertad bajo investigación.
  \item En 2021 Hobby Lobby le demandó en EE.~UU., alegando que le había vendido papiros que no eran suyos.
  \item Obbink ha negado siempre haber obrado mal.
\end{itemize}

No sé con certeza cómo están hoy la causa penal y la demanda civil. Si se necesita para un texto, conviene verificarlo con fuentes actuales.

\section*{Por qué importa para Safo}

Los poemas siguen considerándose \textbf{auténticos}: los textos encajan con lo que se sabe de Safo, y nadie sostiene en serio que sean falsificaciones. Lo que está en entredicho es \textbf{cómo llegaron al mercado}. Probablemente salieron de Egipto de forma ilícita y se presentaron con historias de procedencia inventadas. Por eso muchos papirólogos y algunas revistas consideran éticamente problemático publicar o citar papiros sin procedencia verificable.

El caso dejó a la vista cómo funciona el mercado de antigüedades: la demanda de coleccionistas ricos (aquí, de motivación religiosa) y el prestigio académico, que hace de garantía, se combinan para lavar objetos de origen dudoso.

\end{document}
